\emph{Two-component blocks and the transplanted selection rule.} Let \(X^a=X^a_1\oplus X^a_2\) on \(\mathbb C^N\oplus\mathbb C^N\) be the direct sum of two deformed spheres \(X^a_f=A_f{}^a{}_bJ^b\) of the edge-datum paragraph, and let \(\Theta=\Theta_1\oplus\Theta_2\), with \(\Theta_f\) the spectral projector of \(n_f\cdot J\) on its \(M_f\) largest eigenvalues in summand \(f\), \(2\le M_f<N\); the subsystem of~\cite{FlissFrenkelHartnollSoni2025} attached to \(\Theta\) is invariant under \(U(M_1+M_2)\times U(2N-M_1-M_2)\). Then \(X^a_{\Sigma\bar\Sigma}=c^a_1|u_1)(v_1|+c^a_2|u_2)(v_2|\), with \(u_f,v_f\) in summand \(f\); the nonzero spectrum of \(E=\sum_a\Theta X^a\bar\Theta X^a\Theta\) is \(\{s_1^2,s_2^2\}\), with \(s_f^2=k_fM_f(N-M_f)\) and \(k_f=\tfrac14(\sigma_{f,1}^2+\sigma_{f,2}^2)\), so that \(\operatorname{Tr}E^p=s_1^{2p}+s_2^{2p}\) for every \(p\); and for momenta without entries between the summands the nonzero spectrum of \(Q\) is the union of the two single-cap pairs of the Row length paragraph (exact). The two caps of one axis of the Row length paragraph give the same union, with \(k_1=k_2\), when \(M_1+M_2<N\) and the edges are band-separated; at \(M_1+M_2=N\) their projector is the identity and has no edge. For independent identically distributed Gaussian momenta, which have entries between the summands, the union holds only approximately: at \(M_1=M_2\) and \(s_1=s_2\) the median relative distance of the spectrum from the union of the two pairs formed with the full edge rows, whose norms then also include their entries in the block of the other summand, is \(0.086\) at block rank \(50\) and \(0.023\) at rank \(800\), decaying as about the inverse square root of the rank (numerical). The selection rule of~\cite{Frenkel2023EdgeModes} is derived for the block partition of the matrix degrees of freedom, the cut of the subsystem of~\cite{FlissFrenkelHartnollSoni2025} above, on the fuzzy sphere, in Gaussian states of the quadratic fluctuations about it, and argued for general noncommutative backgrounds: at leading order in \(N\) the moments \(\operatorname{Tr}Q^{2p}\) fix the dominant diagram, with one row for each disconnected component of the cut, of length set by that component, and \(q\) equal rows for each component on the background \(J^a\otimes\mathbf 1_q\). Its use with the supplied states of the Row length paragraph, with rows set by the edge data \(s_f\), and for direct sums of distinct spheres is a transplant (conditional), referred to below as the transplanted selection rule. For a block that is the union of two disconnected components, such as the direct sum above or the two band-separated caps of one axis, whose charge spectrum is exactly the union of the two single-cap pairs, it suggests a two-row diagram of \(U(M_1+M_2)\) with rows fixed by the single-cap edge data; the block \(\Theta_1\otimes\Theta_2\) of the product of two spheres, treated in the toric-block paragraph below, has rank \(M_1M_2\) and a connected edge, so this reading does not apply to it. On the saddle class \(\ell_f=t\,s_f\) with fixed \(t\), the defect \(\Delta=\log\dim_{U(M_1+M_2)}(\ell_1,\ell_2)-\log\dim_{U(M_1)}(\ell_1)-\log\dim_{U(M_2)}(\ell_2)\), \(\ell_1\ge\ell_2\), is \(\sum_{c=1}^{\ell_1}\log\frac{M_1+M_2-1+c}{M_1-1+c}+\sum_{c=1}^{\ell_2}\log\frac{M_1+M_2-2+c}{M_2-1+c}-\log\frac{\ell_1+1}{\ell_1-\ell_2+1}\) (exact, by the hook-content formula; checked in rational arithmetic), the exact form of the excess in item~(ii) of the bearing paragraph (the approximate form holds at leading order for \(\ell_f\ll M_f\)). For the direct sum at \(M_1=M_2=N/2\), \(k_1=k_2=1\), \(t=1\) it grows linearly in \(N\): \(\Delta/N=0.521\) at \(N=3200\), against the limit \(\log\frac{27}{16}=0.5232\) and about \(0.75\) of the approximate form, and \(\sum_{(r,c)\in\mu}M/(M+c-r)=0.811\,|\mu|\) at block rank \(M=M_1+M_2\), so the dependence on the rank survives (numerical). The defect is not monotone in \(\ell_2\): \(\Delta(\ell_1,\ell_2+1)-\Delta(\ell_1,\ell_2)=\log\frac{37}{72}\) at \((M_1,M_2,\ell_1,\ell_2)=(2,24,13,12)\) (exact). Within the dominant-representation class (compact factor, positive contraction, leading semiclassical order with the leading-order entropy of~\cite{FlissFrenkelHartnollSoni2025} and one dominant diagram), with the diagram selected by the transplanted selection rule on the saddle class at fixed frame, the non-additivity of item~(ii) is therefore not an artefact of the fixed probe diagram.

\emph{Row perturbations, stacked rows and the residual bound.} For depths \(1\le q_f\le M_f\) let \(\Delta_{q_1,q_2}\) be the same difference for the diagram of \(U(M_1+M_2)\) with \(q_1\) rows \(\ell_1\) and \(q_2\) rows \(\ell_2\), \(\ell_1\ge\ell_2\ge0\), against \((\ell_1^{q_1})\) of \(U(M_1)\) and \((\ell_2^{q_2})\) of \(U(M_2)\). With \(M=M_1+M_2\), \(\exp\Delta_{q_1,q_2}\) is the product of the content factors \((M+c-r)/(M_1+c-r)\) on the first stack and \((M-q_1+c-r)/(M_2+c-r)\) on the second and of the hook ratios \((\ell_1-c+q_1-r+1)/(\ell_1-c+q_1+q_2-r+1)\) on the first \(\ell_2\) columns of the first stack (exact; checked in rational arithmetic, and equal to the two-row form at \(q_1=q_2=1\)), a function of the ranks and rows alone. With \(M_f=m_fN\), \(u_f=\ell_f/N\), \(G(m,u)=(m+u)\log(m+u)-m\log m-u\) and \(\Lambda=\sum_fq_f[G(m_1+m_2,u_f)-G(m_f,u_f)]\), which is positive and strictly increasing in each \(u_f\), the residual obeys \(|\Delta_{q_1,q_2}-N\Lambda|\le q_1q_2\log(\ell_1+q_1+q_2)+C\), with \(C\) depending only on \(m_f\), \(\ell_f/M\) and \(q_f\) (exact: each row sum is bracketed by its integral, the shift of the ranks by the row index is bounded by the mean-value theorem, and the hook ratios telescope); at equal row fractions \(u_1=u_2\), as at the symmetric saddle, the residual grows as \(q_1q_2\log N\), while at the unequal fractions tested it stays bounded in \(N\) (numerical). So \(\Delta_{q_1,q_2}/N\to\Lambda\), linear in the depths and additive over the stacks: at the symmetric saddle \(\Delta_{q,q}/N=0.5209\), \(1.0379\), \(1.5515\) at \(N=3200\), \(q=1,2,3\), against \(0.5232\,q\) (numerical). The minimum of \(\Delta\) over two-row diagrams with \(|\ell'_f-\ell_f|\le b\) is therefore of order \(N\) for every \(b=o(N)\), with a deficit from the saddle value that is linear in \(b\), of slope \(\sum_f\log\frac{m_1+m_2+u_f}{m_f+u_f}\) at leading order (numerical). Perturbations of order \(N\) remove it: at the symmetric saddle the minimum over \(|\ell'_f-\ell_f|\le cN\) is attained at the far diagonal corner, is positive up to \(c=0.4975\) and vanishes at \(c=\tfrac12\), where the box contains the trivial diagram (numerical, exhaustive at \(N=400\) for \(q_1=q_2\le3\)), whereas removing the second stack leaves \(\Delta_{q,0}/(qN)=0.189\) at \(N=1600\), \(m=(\tfrac12,\tfrac13)\), \(q\le3\) (numerical).

\emph{Frame average at the saddle.} For a summand deformed by \(A_f\), with \(B_f=A_f^{\mathsf T}A_f\), the cap of axis \(n\) has \(s_f^2(n)=\tfrac14M_f(N-M_f)(\operatorname{tr}B_f-n\cdot B_fn)\) (exact; the general-axis form of the edge-datum identity). With \(\rho_f(n)=\tfrac12(\operatorname{tr}B_f-n\cdot B_fn)\), the squared edge-length fraction relative to \(B_f=\mathbf 1\), the saddle rows are \(\ell_f(n)=\ell_f^{(0)}\rho_f(n)^{1/2}\), and \(\Lambda\) depends on the frame through a sum of one term for each summand, strictly increasing and strictly concave in \(\rho_f\). The bound of the preceding paragraph involves the frame only through the rows, which are at most \(\ell_f^{(0)}(\operatorname{tr}B_f/2)^{1/2}\), so the average of \(\Delta_{q_1,q_2}/N\) over the two axes, taken independently and uniformly on the sphere, converges to the sphere mean of \(\Lambda\), uniformly over the frames. When the sphere mean of the squared edge-length fraction equals its round-frame value, \(\operatorname{tr}B_f=3\), that mean is bounded above by the round-frame value, \(0.523248144\,q\) at \(m_f=\tfrac12\), \(\ell_f^{(0)}=N/2\), \(q_1=q_2=q\), with equality only for \(B_f\propto\mathbf 1\), and below by the value at the axis of the largest eigenvalue of \(B_f\), which is positive unless the transverse edge data of a summand degenerate. For \(B_f=\operatorname{diag}(\tfrac12(3-\gamma),\tfrac12(3-\gamma),\gamma)\), at these values and \(q=1\), the averaged limit is \(0.522655\), \(0.519217\), \(0.511350\) and the frame minimum \(0.479485\), \(0.396700\), \(0.268892\) at \(\gamma=1.4\), \(2\), \(2.6\); the finite-\(N\) residuals stay below \(7.7\) for two rows and \(59.8\) at depth three for \(N\le3200\), the stacked limits are the depth times the averaged profile, reproducing the isotropic values of the stacked-rows paragraph at the round frame, and the ratio of \(\mathrm d\log\dim/\mathrm d\log M\) to \(|\mu|\) at the saddle stays between \(0.7858\) and \(0.9023\) (numerical). A sampled frame scan checks the inequality: at \(N=3200\) with symmetric block ranks and saddle parameter one (\(\ell_f^{(0)}=1131\)), over \(150\) sampled frames with \(\operatorname{tr}B=3\) the worst residual is \(7.1\) for two rows and \(54.9\) at depth three, against values of the row-dependent form of the bound between \(8.6\) and \(9.0\) and between \(74.8\) and \(80.2\) (\(8.8\) and \(77.0\) at the round frame), with no violation at \(N=400\), \(1600\) and \(3200\) (numerical, sampled frame grids). This is an average over cap frames at fixed saddle parameter, not the frame average or the minimisation over frames of~\cite{FlissFrenkelHartnollSoni2025}.

\emph{Coincident copies and the multiplicity reading.} The direct sum of \(q\) equal summands, \(\widetilde X=\mathbf 1_q\otimes X\) with \(\widetilde\Theta=\mathbf 1_q\otimes\Theta\) for one cap, is, for undeformed summands, the background \(J^a\otimes\mathbf 1_q\) of the transplanted selection rule; the further conditional reading used here is that multiplicity \(q_f\) of a nonzero eigenvalue of the charge is read as \(q_f\) equal rows, at fixed depths. For copy-invariant momenta \(\mathbf 1_q\otimes\Pi\) the charge is \(\mathbf 1_q\otimes Q\) (exact), so every nonzero eigenvalue has multiplicity \(q\) and the stacked diagram of common depth \(q\) is read; for generic momenta on the same background the rank is the same but the eigenvalues are simple (numerical), so the multiplicity is a property of the state, not of the background alone. For \(q_f\) coincident copies of summand \(f\) of the direct sum, with momenta copy-invariant within each group and without entries between the groups, the charge is \(\bigoplus_f\mathbf 1_{q_f}\otimes Q_f\) (exact), so the stacked diagram \((\ell_1^{q_1},\ell_2^{q_2})\) is read. For the band-separated caps of the Row length paragraph, whose combined spectrum has multiplicity \(1\), no stacked diagram is selected: with saddle rows and depths \((q_1,q_2)\) the stacked moment signature exceeds the multiplicity-\(1\) signature by the fraction \(\sum_f(q_f-1)w_f^p/\sum_fw_f^p\), with \(w_f=k\,m_f(1-m_f)\) and \(k\) the edge-length coefficient of the factor, independent of \(N\) and \(t\): it is \(9/17\) at \(p=1\) and \(81/145\) at \(p=2\) for \((q_1,q_2)=(2,1)\), \((m_1,m_2)=(\tfrac12,\tfrac13)\), \(k=1\), and \(1\) at equal depths \(2\) (exact), a mismatch of leading order in \(N\). For one cap in each of \(q\) summands, or \(q\) coincident copies, the charge is \(-2\mathrm i(UR^\dagger-RU^\dagger)\), with \(U\) the matrix of the orthonormal edge vectors and \(R\) that of the edge rows; with \(C_0=U^\dagger R\), \(Z_0=R^\dagger R\), \(H=(C_0-C_0^\dagger)/(2\mathrm i)\) and \(K=Z_0-C_0^\dagger C_0=R^\dagger(\mathbf 1-UU^\dagger)R\succeq0\), its nonzero spectrum is that of \(-2\mathrm i\left(\begin{smallmatrix}C_0^\dagger&Z_0\\-\mathbf 1_q&-C_0\end{smallmatrix}\right)\), whose characteristic polynomial is \(\det(\lambda^2\mathbf 1_q+4\lambda H-4K)\), so \(\operatorname{Re}C_0\) never enters (exact). If \(H\) and \(K\) are block-scalar, \(h_f\mathbf 1_{q_f}\) and \(n_f\mathbf 1_{q_f}\) with \(n_f>0\), the polynomial is \(\prod_f(\lambda^2+4\lambda h_f-4n_f)^{q_f}\), and the spectrum is the union of \(q_f\) copies of the Row length pair \(-2h_f\pm2(h_f^2+n_f)^{1/2}\); every such datum is realised by \(R=U(A+\mathrm iH)+VK^{1/2}\), with \(A\) real and block-scalar and \(V\) an isometry onto vectors orthogonal to those of \(U\); for \(K=n\mathbf 1_q\) alone the pairs are those of the eigenvalues \(h_f\) of \(H\) (exact). For two diagonal pairs, as for block-diagonal momenta, the quartic \(\lambda^4+a\lambda^3+b\lambda^2+c\lambda+d\) is the square of a real quadratic if and only if \(h_1=h_2\) and \(n_1=n_2\), that is \(\operatorname{Im}\alpha_1=\operatorname{Im}\alpha_2\) and \(|w_{\perp,1}|=|w_{\perp,2}|\), since \(64d-(4b-a^2)^2=-256[(h_1-h_2)^4+2(h_1-h_2)^2(n_1+n_2)+(n_1-n_2)^2]\); a double root need not be a double pair, the pairs \((h,n)=(3,7)\) and \((5,11)\) giving \((\lambda-2)^2(\lambda+14)(\lambda+22)\) (exact). For \(K=n\mathbf 1_q\) a double eigenvalue occurs exactly when two \(h_f\) coincide, a set of real codimension three among Hermitian \(H\), and otherwise the \(2q\) eigenvalues are distinct, with gaps bounded below linearly in \(\min_{f\ne g}|h_f-h_g|\) (exact). At \(h_f=0\) the ratio of \(\sum_je_j^{2p}\) over the nonzero eigenvalues \(e_j\) to its multiplicity-\(1\) value is \(1+\sum_f(q_f-1)n_f^p/\sum_fn_f^p\), the stacked signature above with \(n_f\) in place of \(w_f\), with excess \(9/17\) at \(p=1\) and \(81/145\) at \(p=2\) for \((q_1,q_2)=(2,1)\), \((n_1,n_2)=(\tfrac14,\tfrac29)\) (exact), so in this class the excess is carried by the charge spectrum itself. Within independent Gaussian momenta exact multiplicity is not typical: none of \(6000\) draws at \(q=2\) and \(3\) has it, and along a generic perturbation of copy-invariant momenta the relative gap \(D_{\mathrm{gap}}=\max_{k=1,2}(e_{kq}-e_{(k-1)q+1})/\max_j|e_j|\) of the ordered nonzero eigenvalues grows linearly (numerical). At leading order in \(N\) it is approximately restored when the copies have equal scales: for independent standard complex Gaussian momenta at \(q=2\), with \(M=N/2\), the medians of \(D_{\mathrm{gap}}\) and of \(U_2=\operatorname{Tr}Q^4/\bigl(2q(\operatorname{Tr}Q^2/(2q))^2\bigr)-1\) fall from \(0.28\) and \(0.099\) at \(N=50\) to \(0.11\) and \(0.013\) at \(N=400\), as about \(N^{-1/2}\) and \(N^{-1}\), because the \(h_f\) stay of order one while the \(n_f\) grow as \(M\); with per-copy scales \((1,1.6)\) the median of \(D_{\mathrm{gap}}\) stays near \(0.4\) from \(N=120\) to \(480\) (numerical, sampled edge rows). This is consistent with the \(q\) equal rows of~\cite{Frenkel2023EdgeModes} for equal copies at leading order, and the scale mismatch shows that equal rows at leading order require equal norms of the edge rows.
